\section{Problem Statement}
\label{sec:problem-statement}
Indoor path loss modeling is a critical component of edge IoT design, providing the foundation for network planning, link budget calculations, and data-rate (e.g., \gls{adr}) optimization in complex indoor environments. However, indoor propagation is notoriously difficult to model accurately because radio frequency (\gls{rf}) signals are highly sensitive to building layouts, concrete/brick wall attenuation, and dynamic environmental covariates~\cite{obiri2026environment}. As documented by Obiri and Van Laerhoven~\cite{obiri2026environment}, indoor environments induce rich multipath propagation, sudden attenuation discontinuities, and shadow fading. These factors break the stationarity of propagation channels and yield heavy-tailed, heteroskedastic prediction errors. Consequently, classical log-distance path loss models (\gls{ldplm}) fail to capture this complexity, leading to severely underestimated link uncertainties and inappropriate fade margins.

To solve this, prior work by Obiri and Van Laerhoven~\cite{obiri2026environment,obiri2025comprehensive} proved that integrating co-recorded environmental covariates (such as Relative Humidity, Temperature, and CO\textsubscript{2}) with signal-to-noise ratio (\gls{snr}) metrics significantly improves indoor path loss modeling. Using regularized multiple linear regression and second-order polynomial extensions, they achieved a centralized baseline $R^2 = 0.8219$ (with out-of-sample RMSE of $7.38$\,dB).

However, this centralized approach suffers from four major bottlenecks:
\begin{enumerate}
  \item \textbf{Bandwidth Overhead and Power Depletion (LoRaWAN Constraints):} Streaming continuous, multi-dimensional environmental features (9 variables representing spatial, environmental, and link quality metrics) from edge nodes to a centralized cloud server incurs unsustainable data consumption. Under networks like LoRaWAN, transmitting frequent uplink packets requires substantial time-on-air (\gls{toa}). Because the radio transceiver is the dominant driver of battery drain in low-power microcontrollers (such as the Arduino MKR~WAN~1310), active transmission states quickly deplete the limited battery capacity, reducing device longevity from several years to just a few weeks. Furthermore, this high transmission frequency violates regional wireless regulations, specifically the strict 1\,\% duty cycle limit in the European 868\,MHz band, resulting in severe channel congestion, high packet collision rates, and network capacity degradation.
  \item \textbf{Data Privacy Vulnerabilities:} Environmental covariates like CO\textsubscript{2} and ambient temperature act as direct proxies for human activity and building occupancy (revealing exactly when rooms are occupied and how many people are present). Streaming these raw, high-resolution telemetry logs to a centralized server exposes private indoor dynamics and employee routines, violating strict data privacy regulations such as the GDPR.
  \item \textbf{Static Deployment and Adaptation Barrier:} A centralized model is static and trained on historical dataset conditions. If physical modifications occur in the indoor environment (such as adding walls, relocating furniture, or moving sensor nodes to new rooms), the propagation characteristics change, and the centralized model's accuracy permanently degrades. Adapting the model requires initiating another expensive centralized data-collection and retraining cycle.
  \item \textbf{Fade Margin Miscalibration under Non-Stationarity:} Indoor environments induce non-stationarity, meaning link characteristics fluctuate dynamically over time. If a static centralized model cannot adapt on the fly, its residual prediction errors drift, leading to miscalibrated fade margins. This results in either underestimated fade margins (causing connection dropouts and low link reliability) or overly conservative margins (causing excessive transmission power consumption and unnecessary energy waste).
\end{enumerate}

While moving to a decentralized Federated Learning (\gls{fl}) setup keeps data local and solves the privacy and bandwidth bottlenecks, it introduces the core machine learning challenge of this thesis: \emph{Client Drift}~\cite{ikram2026empirical}. Because the LoRaWAN nodes are deployed in physically distinct rooms (e.g., stairwells, kitchens, offices), their local data distributions are highly heterogeneous, or non-\gls{iid} (Independent and Identically Distributed). When standard federated algorithms like Federated Averaging (\gls{fedavg}) are applied, the local models optimize solely for their specific local environments and diverge. When these drifted weights are averaged on the server, the resulting global model suffers from severe accuracy degradation, failing to generalize across the building. Thus, the core challenge is designing a decentralized model that can be trained under strict memory limits (TinyML) and communication bounds (LoRaWAN) without succumbing to non-IID client drift.
